\documentclass[11pt]{article}
\usepackage[utf8]{inputenc}
\usepackage{amsmath,amssymb}
\usepackage[margin=1in]{geometry}
\usepackage{hyperref}
\emergencystretch=3em

\title{The reduced form of the two-dimensional cutoff expansion}
\author{Claude, with Daniel Murfet}
\date{2026-09-07}

\begin{document}
\maketitle
\sloppy

\begin{abstract}
The cutoff theorem of unit~107 lists the transferred terms of all faces and corners. The monomial
terms cancel exactly against the corner terms (Astra~\#5(a)): the $m=j$ monomial coefficient of the
$u$-face $i$ equals the finite-part coefficient of the corner $(i,j)$, and the $m=i$ monomial
coefficient of the $v$-face $j$ equals the monomial coefficient of the corner, as functions of $s$
and with equal exponents, both at resonance and off it. What remains is the reduced sum
\begin{gather*}
\sum_{i<M_1}N^{-(p+i/k_1)}\mathcal M\bigl[k_1^{-1}\mathrm{FP}_{k_2i/k_1}(a_i)\bigr]
+\sum_{j<M_2}N^{-(p+j/k_2)}\mathcal M\bigl[k_2^{-1}\mathrm{FP}_{k_1j/k_2}(b_j)\bigr]\\
+(\text{$u$-face logs})+(\text{$v$-face logs})-(\text{corner logs}),
\end{gather*}
and \texttt{twoD\_cutoff\_reduced\_isBigO} restates unit~107 with it: the constant coefficients of
the higher-order $d=2$ expansion are exactly the weighted finite parts of the face jets, as
predicted. The three log families combine to one copy per collision; that final simplification is
left as a separate lemma. Zero \texttt{sorry}/\texttt{axiom}.
\end{abstract}

\section{The things formalised}
{\small
\begin{verbatim}
theorem cornerData_sum (c : ℝ → ℝ) (g : ℕ → ℝ) (M : ℕ) (hM : 0 < M) (s : ℝ) :
    ∑ m ∈ Finset.range M, cornerData c m s * g m = c s * g 0
theorem faceCoeff_cornerData_succ (γ A : ℝ) (k₁ k₂ : ℕ) (c : ℝ → ℝ) (m : ℕ) :
    faceCoeff γ A k₁ k₂ (cornerData c) (m + 1) = fun _ => 0
theorem faceFPCoeff_corner (γ b : ℝ) (k₂ : ℕ) (c : ℝ → ℝ) (M : ℕ) (hM : 0 < M) (s : ℝ) :
    faceFPCoeff γ b k₂ (fun _ s => c s) (cornerData c) M s
      = 1 / (k₂ : ℝ) * (c s * axisPrim γ b 0)
theorem uface_coeff_eq_corner_fp (b p : ℝ) (h₁ h₂ k₁ k₂ i j Mc : ℕ) (hb : 0 < b)
    (hk₁ : 0 < k₁) (hk₂ : 0 < k₂) (hp₁ : ((h₁ : ℝ) + 1) / k₁ = p)
    (hp₂ : ((h₂ : ℝ) + 1) / k₂ = p) (hMc : 0 < Mc) (c : ℕ → ℕ → ℝ → ℝ) :
    faceCoeff ((k₂ : ℝ) * ((((h₁ + i : ℕ) : ℝ) + 1) / k₁) - h₂ - 1) (b ^ k₁) k₂ k₁
        (fun m => c i m) j
      = faceFPCoeff ((k₁ : ℝ) * ((((h₂ + j : ℕ) : ℝ) + 1) / k₂) - ((h₁ + i : ℕ) : ℝ) - 1) b k₂
        (fun _ s => c i j s) (cornerData (c i j)) Mc
theorem vface_coeff_eq_corner_coeff (b : ℝ) (h₁ h₂ k₁ k₂ i j : ℕ) (c : ℕ → ℕ → ℝ → ℝ) :
    faceCoeff ((k₁ : ℝ) * ((((h₂ + j : ℕ) : ℝ) + 1) / k₂) - h₁ - 1) (b ^ k₂) k₁ k₂
        (fun m => c m j) i
      = faceCoeff ((k₁ : ℝ) * ((((h₂ + j : ℕ) : ℝ) + 1) / k₂) - ((h₁ + i : ℕ) : ℝ) - 1)
        (b ^ k₂) k₁ k₂ (cornerData (c i j)) 0
theorem uface_exp_eq / vface_exp_eq   -- the matching exponents
noncomputable def reducedSum (β b : ℝ) (h₁ h₂ k₁ k₂ M₁ M₂ : ℕ) (a bj : ℕ → ℝ → ℝ → ℝ)
    (c : ℕ → ℕ → ℝ → ℝ) (N : ℝ) : ℝ   -- face finite parts + u-logs + v-logs − corner logs
theorem overcomplete_eq_reduced ... :
    (∑ i ∈ Finset.range M₁, faceSum [u-face i] N) + (∑ j ∈ Finset.range M₂, faceSum [v-face j] N)
      - ∑ i ∈ Finset.range M₁, ∑ j ∈ Finset.range M₂, faceSum [corner i j] N
      = reducedSum β b h₁ h₂ k₁ k₂ M₁ M₂ a bj c N
theorem twoD_cutoff_reduced_isBigO ... :
    (fun N : ℝ => twoDAmp β b N h₁ h₂ k₁ k₂ Φ - reducedSum β b h₁ h₂ k₁ k₂ M₁ M₂ a bj c N)
      =O[atTop] fun N : ℝ =>
        N ^ (-(p + min ((M₁ : ℝ) / k₁) ((M₂ : ℝ) / k₂))) * (1 + Real.log N)
\end{verbatim}
}

\section{Mathematics}

With $\gamma_u=k_2i/k_1$ (the $u$-face) and $\gamma_c=k_1j/k_2-i$ (the corner), $j=\gamma_u$ iff
$\gamma_c=0$ iff $k_1j=k_2i$. Off resonance the $u$-face monomial coefficient is
$-k_1^{-1}(b^{k_1})^{-(j-\gamma_u)/k_2}c_{ij}/(j-\gamma_u)=-b^{-\gamma_c}c_{ij}/(k_2\gamma_c)$,
which is $k_2^{-1}c_{ij}\,\mathrm{axisPrim}(\gamma_c,b,0)$, the finite part of the corner (whose
regularised integral vanishes since its Taylor remainder is identically zero); at resonance both are
$k_2^{-1}c_{ij}\log b$ (using $\log b^{k_1}=k_1\log b$). The $v$-face/corner monomial identity is a
relabelling $\gamma_c=\gamma_v-i$. The corner's monomial terms with $m\ge1$ have zero coefficient
functions and so vanish under the moment functional. Summing over faces and corners and commuting
the double sum gives the reduced form; no integrability or asymptotics is used, only pointwise
equality of coefficient functions.

\section{Paper-facing meaning}

This is the grouped coefficient description of Astra~\#5(a) up to the merging of the three log
families: at exponent $p+q$ the nonlogarithmic coefficient is $\mathcal M_{p+q}$ of
$k_1^{-1}\mathrm{FP}_{k_2q}(a_i)+k_2^{-1}\mathrm{FP}_{k_1q}(b_j)$ (whichever of $q=i/k_1$,
$q=j/k_2$ apply), and the corner monomials contribute nothing beyond what is already inside the two
face finite parts. This is the closest formal counterpart so far of the paper's description of the
higher-order coefficients through renormalised face integrals.

\section{Lean notes}

\begin{itemize}
\item Two competing normalisations of the same denominator ($k_1j-k_2i$ vs $-(k_1j)+k_2i$) defeat
\texttt{field\_simp}; rewrite both denominators to a common form first
(\texttt{d1}, \texttt{d2}) and then \texttt{field\_simp} closes the goal outright.
\item \texttt{Fin.sum\_univ\_eq\_sum\_range} with a large explicit lambda left an instance stuck
inside \texttt{rw}; \texttt{congr 1} down to the sum and \texttt{rw [Finset.sum\_range]} on the
range side, then \texttt{Finset.sum\_congr} over \texttt{Fin}, is robust. \texttt{Fin.pos j} gives
\texttt{0 < M₂} from an inhabitant.
\item \texttt{Finset.sum\_comm} as a bare \texttt{rw} may hit the wrong double sum; state the
specific commutation as a \texttt{have} with explicit summands.
\item \texttt{try ring} after \texttt{field\_simp} accommodates both outcomes (closed or not).
\end{itemize}

\end{document}
